\section{Full derivation of the $d\ge4$ exclusions}

This chapter gives the complete derivation of the $d\ge4$ exclusions.
The proof splits according to the location and weight of the first
big step.

\subsection{The three branch types}

For $d\ge4$, the first big step of the candidate block is one of
three types:
\begin{enumerate}[leftmargin=*]
\item the incoming step $g_{\mathrm{prev}}\to g$, which happens when
  $e\ge3$;
\item the step $y\to x$, which happens when $e=2$ and $a\ge1$ and has
  weight $1+4a\ge5$;
\item the step $z\to w$, which happens when $e=2$ and $a=0$ and has
  weight $5$ or $6$.
\end{enumerate}

\subsection{The finite base}

The finite base says:
\begin{itemize}[leftmargin=*]
\item for every $n<15$, $t_n\le2$;
\item $t_{15}=3$;
\item the first big step is unique.
\end{itemize}
These facts are compiled in
\texttt{fullOrbit\_\allowbreak first\_\allowbreak t\_\allowbreak ge3\_
\allowbreak is\_\allowbreak exactly\_\allowbreak 3} and
\texttt{first\_\allowbreak big\_\allowbreak step\_\allowbreak unique}.

\subsection{Branch $e\ge3$}

The incoming step $g_{\mathrm{prev}}\to g$ has weight $e\ge3$, so it
is a big step.  The finite base forces
\begin{equation*}
  e=3,\qquad j=17.
\end{equation*}
For this branch the candidate residue is then excluded modulo $640$
or $1280$.  The two compiled statements are
\texttt{dge4\_\allowbreak e3\_\allowbreak j17\_\allowbreak t1\_\allowbreak
excluded} and
\texttt{dge4\_\allowbreak e3\_\allowbreak j17\_\allowbreak
t2\_\allowbreak delta3\_\allowbreak excluded}.

\subsection{Branch $e=2$, $a\ge1$}

The step $y\to x$ has weight
\begin{equation*}
  1+4a\ge5.
\end{equation*}
This would be the first big step of the orbit prefix.  But the finite
base says the first big step has weight exactly $3$, so this branch is
impossible.  The compiled statement is
\texttt{dge4\_\allowbreak e2\_\allowbreak a\_\allowbreak ge1\_\allowbreak
excluded}.

\subsection{Branch $e=2$, $a=0$}

The step $z\to w$ has weight $5$ or $6$, again contradicting the
finite base.  The same uniqueness argument applies: the true first big
step occurs at depth $15$ with weight $3$, so a later candidate big
step cannot be first, and an earlier one cannot have weight $5$ or
$6$.

\subsection{Why uniqueness is used}

The uniqueness statement says that any depth whose earlier weights are
all at most $2$ and whose own weight is at least $3$ must be depth
$15$ with weight $3$.  A candidate first big step of weight at least
$5$ fails this conclusion, regardless of where it is placed.

\subsection{Compiled statements}

\begin{itemize}[leftmargin=*]
\item \texttt{dge4\_\allowbreak e2\_\allowbreak a\_\allowbreak ge1\_
\allowbreak excluded};
\item \texttt{dge4\_\allowbreak e3\_\allowbreak j17\_\allowbreak
t1\_\allowbreak excluded};
\item \texttt{dge4\_\allowbreak e3\_\allowbreak j17\_\allowbreak
t2\_\allowbreak delta3\_\allowbreak excluded};
\item \texttt{candidate\_\allowbreak first\_\allowbreak big\_\allowbreak
step\_\allowbreak weight\_\allowbreak ge\_\allowbreak five}.
\end{itemize}

\subsection{Status}

\begin{center}
\small
\begin{tabular}{@{}p{0.56\textwidth}p{0.36\textwidth}@{}}
\toprule
Item & Status \\
\midrule
branch classification & math proven \\
$e=2$ branches & math proven, Lean compiled \\
$e=3$, $j=17$ residues & math proven, Lean compiled \\
\bottomrule
\end{tabular}
\end{center}
